% TOP_PROBLEM: {"id":30007734,"problem_number":"LOCAL-30007734","title":"Effective uses of school funding","category_id":21,"category":{"id":21,"name":"economics","display_name":"Economics","description":"Open economic theory statements, published research agendas, and proposed scoped specifications; question provenance and historical priority are qualified per record.","slug":"economics","order_index":21,"created_at":"2026-10-08T00:00:00Z"},"status":"research_question","external_url":null,"legacy_ids":[],"published":false,"statement":"Compare feasible uses of additional school funds at a fixed budget and determine whether allocation rankings transport to another defined district population.","economics":{"original_id":223,"field":"Education, families and demography","kind":"design","formalization_basis":"adapted_research_question","source_status":"research_agenda","priority_status":"earliest_found","priority_note":"Earliest explicit statement located in this audit; no exhaustive priority search or first-ever claim. The mathematical specification below is adapted, not quoted from the source.","earliest_verified_reference":{"authors":["Danielle V. Handel","Eric A. Hanushek"],"title":"Contexts of Convenience: Generalizing from Published Evaluations of School Finance Policies","year":2024,"url":"https://www.hanushek.net/publications/contexts-convenience-generalizing-published-evaluations-school-finance-policies-1","doi":"10.1177/0193841X241228335","locator":"Abstract on the author’s publication page","evidence_summary":"The authors state that existing evaluations do not adequately describe effective contexts or how extra funds are used."},"current_reference":{"authors":["Danielle V. Handel","Eric A. Hanushek"],"title":"Contexts of Convenience: Generalizing from Published Evaluations of School Finance Policies","year":2024,"url":"https://www.hanushek.net/publications/contexts-convenience-generalizing-published-evaluations-school-finance-policies-1","doi":"10.1177/0193841X241228335","locator":"Abstract on the author’s publication page","evidence_summary":"The authors state that existing evaluations do not adequately describe effective contexts or how extra funds are used."},"formalization_note":"The source explicitly identifies the underlying gap or agenda. Population, horizon, interventions, estimands, and auxiliary assumptions are proposed operational choices, not a canonical source theorem. Partial later answers are compatible with this scoped question; the citation alone does not prove absence of every subsequent solution. A 2023 working-paper predecessor exists, but its explicit passage was not inspected; the verified gap statement is in the 2024 author publication abstract. Mathematical scope was checked independently; added definedness, feasibility, or identification conditions belong to this proposed formalization."},"statusQualification":"Published question or research agenda verified at the cited locator. The mathematical specification is adapted for this catalogue and includes additional assumptions; source publication does not establish first-ever priority or certify this exact model remains unsolved. The source explicitly identifies the underlying gap or agenda. Population, horizon, interventions, estimands, and auxiliary assumptions are proposed operational choices, not a canonical source theorem. Partial later answers are compatible with this scoped question; the citation alone does not prove absence of every subsequent solution. A 2023 working-paper predecessor exists, but its explicit passage was not inspected; the verified gap statement is in the 2024 author publication abstract. Mathematical scope was checked independently; added definedness, feasibility, or identification conditions belong to this proposed formalization.","statusReviewedAt":"2026-10-08"}
% FIRST_POSTED: 2026-10-08
% ENTRY_KIND: statement_only
% !TeX program = lualatex
\documentclass[11pt]{article}
\usepackage{fontspec}
\usepackage[a4paper,margin=25mm]{geometry}
\usepackage{amsmath,amssymb,mathtools}
\usepackage{xurl}
\usepackage[hidelinks]{hyperref}
\newcommand{\sref}[2]{\hyperlink{source:#1:#2}{[S#2]}}
\setlength{\emergencystretch}{3em}
\begin{document}
% problemId:30007734
\section*{Effective uses of school funding}\label{30007734}
\subsection{Definitions and mathematical statement}
Consider public schools serving disadvantaged pupils in one state over five budget years. Let \(b>0\) be additional annual funding per pupil and \(x=(x_1,\ldots,x_J)\) a spending allocation across \(J\) prespecified uses, such as tutoring, teacher retention, class-size reduction, and facilities. Feasibility requires \(x_j\geq0\) and \(\sum_{j=1}^{J}x_j=b\). Let \(Y_i(x)\) denote pupil \(i\)'s educational attainment and independently scaled learning outcomes, combined using fixed weights announced before estimation. Define
\[x^*(b)\in\arg\max_{x\geq0:\sum_jx_j=b}E[Y_i(x)-Y_i(0)].\]
Assume finite expected outcomes and a continuous mean response on the compact allocation simplex, or report a supremum if attainment cannot be justified. Estimate credible response surfaces or feasible policy rankings using randomized spending packages or justified administrative variation. School context, staff availability, implementation quality, and baseline resources must be measured; average funding effects do not identify this allocation problem. Address substitution of local funds and spillovers from teacher recruitment. Report whether rankings transport to a separately specified district population, with overlap conditions and sensitivity to unobserved effect modifiers. The source explicitly says existing evaluations inadequately explain effective contexts and uses of funds. The budget, allocation set, outcome index, and transport population are proposed details rather than a source-given canonical optimization problem.

\par\textbf{Formulation and scope.} The source explicitly identifies the underlying gap or agenda. Population, horizon, interventions, estimands, and auxiliary assumptions are proposed operational choices, not a canonical source theorem. Partial later answers are compatible with this scoped question; the citation alone does not prove absence of every subsequent solution. A 2023 working-paper predecessor exists, but its explicit passage was not inspected; the verified gap statement is in the 2024 author publication abstract. Mathematical scope was checked independently; added definedness, feasibility, or identification conditions belong to this proposed formalization.

\par\textbf{Open-question provenance.} An explicit question or research agenda was verified at the source locator. This does not by itself certify that every later version remains unresolved \sref{30007734}{1}.

\par\textbf{Historical priority.} Earliest explicit statement located in this audit; no exhaustive priority search or first-ever claim. The mathematical specification below is adapted, not quoted from the source.

\subsection{Short English statement}
Compare feasible uses of additional school funds at a fixed budget and determine whether allocation rankings transport to another defined district population.

\subsection{Sources}
\begin{itemize}\raggedright
\item[\textnormal{[S1]}]\hypertarget{source:30007734:1}{} Danielle V. Handel, Eric A. Hanushek. 2024. Contexts of Convenience: Generalizing from Published Evaluations of School Finance Policies. Abstract on the author’s publication page.
\url{https://www.hanushek.net/publications/contexts-convenience-generalizing-published-evaluations-school-finance-policies-1}.
DOI: 10.1177/0193841X241228335.
{\small\textit{Earliest verified question/agenda reference; historical priority qualified below. The authors state that existing evaluations do not adequately describe effective contexts or how extra funds are used.}}
\end{itemize}
\end{document}
